\subsection{initramfs}
\label{subsec:initramfs}

这是启动时所加载的文件系统，以 cpio 格式打包。我生成的镜像含有一个名为 /init 的 ELF 文件，以及三个文件夹 /tmp, /dev 和 /mnt 用于后续的 mounting。

这里的 cpio 用的是 newc 格式\cite{cpionewc}：每个字段都是一个字符数组，所以甚至可以拿文本编辑器打开阅读。这也避免了端序的问题。解析起 cpio 来也十分简单，甚至比 FDT 的格式还要更加明了一些，只是由于路径是扁平的，我们需要合适地分割字符串来得到目录名称和文件名称，所以一个趁手的 string 类是很有用的。这也是我在刚启动时就几乎立即激活内存分配器的原因之一。

这个 initramfs 是只读的，而且在整个 OS 的生命中都不会被释放：毕竟 /init 一直在跑。不过它除了最开始启动时的 mount 以外，后续就几乎不会再使用了。

目前我还没有支持压缩过后的镜像，例如 .cpio.gz 之类的。实际上这个 cpio 只有几百个字节，因此我觉得压缩是不必要的。基于同样的理由，有很多“看来用不上”的内容我都没有在这个文件系统里实现，例如 symlink。不过它们在\cref{subsec:tmpfs}处所描述的 tmpfs 中都有实现。
